% minilecture03.tex: Tutorial 3 mini-lecture: Atmospheres
% Planetary Systems, Kapteyn Institute, University of Groningen
% A 5-10 minute recap, presented by the TAs at the start of Tutorial 3,
% of the Lecture 5-6 concepts that Worksheet 3 trains.

\documentclass[aspectratio=169, 11pt]{beamer}

% ── Theme and macros (shared with the lecture decks) ────────
\usepackage{../../slides/common/beamerthemeIPS}
% macros.tex: Shared math macros for IPS lecture slides
% Mirrors definitions in book/_config.yml for consistency
% Uses \providecommand to avoid conflicts with unicode-math

% ── Derivatives ─────────────────────────────────────────────
\providecommand{\dv}[2]{\frac{\mathrm{d} #1}{\mathrm{d} #2}}
\providecommand{\dd}{}\renewcommand{\dd}{\mathrm{d}}
\providecommand{\pdv}[2]{\frac{\partial #1}{\partial #2}}

% ── Solar quantities ────────────────────────────────────────
\newcommand{\Msun}{M_{\odot}}
\newcommand{\Rsun}{R_{\odot}}
\newcommand{\Lsun}{L_{\odot}}

% ── Planetary quantities ────────────────────────────────────
\newcommand{\Mearth}{M_{\oplus}}
\newcommand{\Rearth}{R_{\oplus}}
\newcommand{\Mjup}{M_{\mathrm{J}}}
\newcommand{\Rjup}{R_{\mathrm{J}}}

% ── Constants ───────────────────────────────────────────────
\newcommand{\kB}{k_{\mathrm{B}}}

% ── Media links ─────────────────────────────────────────────
% Clickable button that opens the animated or video version of a
% figure, e.g. \mediabutton{https://...gif}{Play animation}
\providecommand{\mediabutton}[2]{\href{#1}{\beamergotobutton{#2}}}

\graphicspath{{figures/}{../../slides/lecture06/figures/}{../../slides/lecture05/figures/}{../../slides/common/}}

% ── Metadata ────────────────────────────────────────────────
\title[T3 mini-lecture: Atmospheres]{Tutorial 3 Mini-Lecture:\\Atmospheres}
\subtitle{The concepts behind Worksheet 3 \textbullet\ Lectures 5--6}
\author{}
\institute{Kapteyn Astronomical Institute\\University of Groningen}
\date{2026}
\titlebgcredit{Jupiter's atmosphere. Credit: NASA/JPL-Caltech/SwRI/MSSS}
\titlebgopacity{0.55}

% ════════════════════════════════════════════════════════════
\begin{document}

% ── Title slide ─────────────────────────────────────────────
\begin{frame}[plain,noframenumbering]
  \titlepage
\end{frame}

% ── Roadmap ─────────────────────────────────────────────────
\begin{frame}{What Worksheet 3 trains}
  Four problems, all built on Lectures 5--6. The physics you need, per problem:
  \vskip6pt\relax
  \begin{enumerate}
    \item \textbf{How thick is an atmosphere?}: scale height, hydrostatic balance, column mass.
    \item \textbf{Sunlight in, infrared out}: equilibrium temperature, the greenhouse effect.
    \item \textbf{From vapour to cloud base}: Clausius-Clapeyron, dew point, the LCL, the saturation-curve figure.
    \item \textbf{Climate on the edge}: the faint young Sun, ice-albedo bistability, the carbonate-silicate thermostat.
  \end{enumerate}
  \vskip2pt\relax
  \keyresult{Every problem has the shape and level of one exam question. All parts are analytical; no computer is needed.}
\end{frame}

% ── P1 ──────────────────────────────────────────────────────
\begin{frame}{How thick is an atmosphere?}
  \small
  \begin{columns}[T]
    \column{0.52\textwidth}
    Two workhorses of Lecture 5:
    \[
      \boxed{H = \frac{\kB T}{\mu\,m_u\,g}}
      \qquad
      P_0 = m_{\mathrm{col}}\,g
    \]
    \begin{itemize}
      \item $H$ compares thermal energy against gravitational confinement: Earth 8.4 km, Titan 21 km.
      \item Integrating $\dd P = -\rho g\,\dd z$ over the column: surface pressure is the weight of the air overhead.
    \end{itemize}

    \column{0.44\textwidth}
    \begin{block}{Watch the ratio, not the terms}
      $H \propto T/g$: when two effects pull in opposite directions, put numbers on both ratios before deciding which one wins. Problem 1(b) turns exactly on this.
    \end{block}
  \end{columns}
  \vskip2pt\relax
  \keyresult{Ten tonnes of air stand on every square metre of Earth; the scale height says how thickly it is spread.}
\end{frame}

% ── P2 ──────────────────────────────────────────────────────
\begin{frame}{Sunlight in, infrared out}
  \small
  \begin{columns}[T]
    \column{0.52\textwidth}
    Balance absorbed sunlight against emitted infrared:
    \[
      \boxed{(1-A)\,\frac{S}{4} = \sigma T_{\mathrm{eq}}^4}
    \]
    \begin{itemize}
      \item $S$ falls with the inverse square of distance; $A$ is the Bond albedo.
      \item Venus: brilliant clouds give $T_{\mathrm{eq}} = 227$ K, colder than Earth's 255 K, above a 737 K surface.
    \end{itemize}

    \column{0.44\textwidth}
    \begin{block}{The greenhouse in one sentence}
      The atmosphere passes sunlight but blocks the surface's infrared, so the surface must heat up until its share leaks through.
    \end{block}
  \end{columns}
  \vskip2pt\relax
  \keyresult{The place to test any claim about the greenhouse budget is the top of the atmosphere: check whether absorbed sunlight and outgoing infrared still balance there.}
\end{frame}

% ── P3 ──────────────────────────────────────────────────────
\begin{frame}{From vapour to cloud base}
  \small
  \begin{columns}[T]
    \column{0.52\textwidth}
    The blackboard result of Lecture 6:
    \[
      \boxed{P_{\mathrm{sat}}(T) = P_{\mathrm{ref}}\,
      e^{-\frac{L_v}{R_v}\left(\frac{1}{T} - \frac{1}{T_{\mathrm{ref}}}\right)}}
    \]
    \begin{itemize}
      \item Invert it for the dew point $T_d$; rising air cools at $\Gamma_d = 9.8$ K\,km$^{-1}$ until it gets there.
      \item $z_{\mathrm{LCL}} = (T - T_d)/\Gamma_d$: the cloud base. Problem 3(a) runs this chain from surface humidity to cloud-base height.
    \end{itemize}

    \column{0.44\textwidth}
    \centering
    \includegraphics[width=\linewidth, height=0.42\textheight, keepaspectratio]{psat_curves}\\
    {\tiny The saturation curves printed on the sheet: Problem 3(b) finds the species whose curve passes the given $(T, P)$ point, the body follows from the given data. Argue the species from the curves.}
  \end{columns}
  \vskip2pt\relax
  \keyresult{One exponential governs every cloud in the solar system; $L_v/R_v$ sets how sharply it bites.}
\end{frame}

% ── P4 ──────────────────────────────────────────────────────
\begin{frame}{Climate on the edge}
  \small
  \begin{columns}[T]
    \column{0.52\textwidth}
    Long-term climate is a balance with feedbacks:
    \begin{itemize}
      \item Young Sun at 71\%: $T_{\mathrm{eq}} = 234$ K, yet the rocks record liquid water: the faint young Sun paradox.
      \item Ice-albedo feedback: two stable climates under the same Sun.
      \item Carbonate-silicate thermostat: warmer $\rightarrow$ faster weathering $\rightarrow$ less $\mathrm{CO_2}$ $\rightarrow$ cooler.
    \end{itemize}

    \column{0.44\textwidth}
    \centering
    \includegraphics[width=\linewidth, height=0.42\textheight, keepaspectratio]{snowball_bistability_unlabeled}\\
    {\tiny The bistability diagram printed on the sheet: Problem 4(b) classifies each crossing from the local slopes, by asking whether absorption or emission grows faster after a small nudge in temperature.}
  \end{columns}
  \vskip2pt\relax
  \keyresult{A thermostat needs a working source and a working sink: Venus lost its water, Mars lost its volcanism, Earth kept both.}
\end{frame}

% ── Closer ──────────────────────────────────────────────────
\begin{frame}{Working the worksheet}
  \begin{itemize}
    \item \textbf{Show your algebra.} The exam asks for the steps, not only the number.
    \item Carry \textbf{4 significant figures} through intermediate steps; round at the end.
    \item Use the \textbf{checkpoints} (``show that \ldots'') to verify your work and to keep moving if a step resists.
    \item Qualitative parts want \textbf{2 to 3 sentences} of physics, not an essay.
    \item Two parts read figures printed on the sheet (the saturation curves and the bistability diagram); argue the species and the climate branch from the curves.
  \end{itemize}
  \vskip8pt\relax
  \keyresult{Worksheet and full solutions: \href{https://ips.formingworlds.space/worksheets.html}{\texttt{ips.formingworlds.space/worksheets.html}}. We discuss questions, not presentations of the solutions.}
\end{frame}

\end{document}
